\documentclass{article}
\usepackage[UTF8]{ctex}
\usepackage{amsmath,amssymb,amsfonts}
\usepackage{enumitem}
\usepackage{geometry}
\geometry{a4paper, margin=1in}
% === 错题本排版优化（不参与题目标识） ===
\setlength{\parindent}{0pt}
\setlength{\parskip}{0.35em plus 0.10em minus 0.08em}
\setlength{\emergencystretch}{3em}
\setlist{itemsep=0.25em, topsep=0.40em, parsep=0pt}
\allowdisplaybreaks[2]
% === 错题本排版优化结束 ===
\begin{document}
\section*{长安大学2025年高等代数真题}
\section*{175．长安大学2025年高等代数试卷}
一．填空题．每题5分，共30分．\\

\subsection{题目1}
1．多项式 \(f(x)=x^{3}-3 x+k\) 有重根的充要条件为 \(\_\_\_\_\)．

\subsection{题目2}
2．设 4 阶非零方阵 \(A=\left(a_{i j}\right)\) 满足 \(A_{i j}=a_{i j}\) ，则 \(|A|=\) \(\_\_\_\_\)．

\subsection{题目3}
3．设 \(A=\left(\begin{array}{ccc}2 & -1 & 1 \\ x & 5 & -2 \\ -3 & -3 & b\end{array}\right)\) 可对角化且有二重特征值，则 \(x=\) \(\_\_\_\_\).\\
\(\_\_\_\_\)

\subsection{题目4}
4．设 \(A=\left(\begin{array}{rrr}3 & 0 & 0 \\ 0 & -1 & 4 \\ -1 & -1 & 3\end{array}\right)\) ，则 \(A\) 的若尔当标准形为 \(\_\_\_\_\)．

\subsection{题目5}
5．设 \(A\) 为 \(n\) 阶退化矩阵，存在 \(A_{i j} \neq 0\) ，则 \(A^{*}\) 的秩为 \(\_\_\_\_\)．\\
\(\_\_\_\_\)

\subsection{题目6}
6．在线性空间 \(\mathbb{R}^{2 \times 2}\) 中，当 \(x=\) \(\_\_\_\_\)时，向量

\[
A_{1}=\left(\begin{array}{ll}
x & 1 \\
1 & 1
\end{array}\right), A_{2}=\left(\begin{array}{ll}
1 & x \\
1 & 1
\end{array}\right), A_{3}=\left(\begin{array}{ll}
1 & 1 \\
1 & x
\end{array}\right), A_{4}=\left(\begin{array}{ll}
1 & 1 \\
x & 1
\end{array}\right)
\]

线性相关。

二．（10 分）求多项式 \(f(x)=x^{4}-6 x^{3}+10 x^{2}-2 x-3\) 在有理数域 \(\mathbb{Q}\) 及实数域 \(\mathbb{R}\) 上的标准分解式．

三．（15 分）叙述并证明艾森斯坦判别法．

四．（10 分）计算行列式

\[
D_{n}=\left|\begin{array}{cccccc}
3 & 2 & 0 & \cdots & 0 & 0 \\
1 & 3 & 2 & \cdots & 0 & 0 \\
0 & 1 & 3 & \cdots & 0 & 0 \\
\vdots & \vdots & \vdots & & \vdots & \vdots \\
0 & 0 & 0 & \cdots & 3 & 2 \\
0 & 0 & 0 & \cdots & 1 & 3
\end{array}\right| .
\]

五．（10分）设 \(A, B\) 为 \(n\) 阶实方阵，满足 \(A^{2}=B^{2}=E\) ，且 \(|A|+|B|=0\) ，证明：\(|A+B|=0\) 。

六．（10分）设 \(A\) 为 \(n\) 阶方阵，证明：\(R\left(A^{n}\right)=R\left(A^{n+1}\right)\) ．

七．（15 分）设 \(A\) 为实反对称矩阵．证明：\(E-A^{2}\) 非奇异．

八．（15 分）设 3 阶实对称矩阵 \(A\) 的特征值为 \(\lambda_{1}=-1, \lambda_{2}=\lambda_{3}=1, \xi_{1}=(0,1,1)^{\mathrm{T}}\) 是 \(A\) 对应于 \(\lambda_{1}\) 的一个特征向量，求 \(A^{2025}\) 。

九．（20 分）设 \(V=\{f(x) \in P[x] \mid \operatorname{deg} f(x)<n\) 或 \(f(x)=0\}\) ，作 \(V\) 的变换 \(\sigma\) 如下：对任意的\(f(x) \in V\) ，令 \(\sigma(f(x))=x f^{\prime}(x)-f(x)\) ．

\begin{itemize}

  \item[1．] 证明：\(\sigma\) 是 \(V\) 上的线性变换．

  \item[2．] 求 \(\operatorname{Ker} \sigma\) 及 \(\operatorname{Im} \sigma\) ．

  \item[3．] 证明：\(V=\operatorname{Ker} \sigma \oplus \operatorname{Im} \sigma\) ．
\end{itemize}

十．（15 分）已知二次型

\[
f\left(x_{1}, x_{2}, x_{3}\right)=5 x_{1}^{2}+5 x_{2}^{2}+c x_{3}^{2}-2 x_{1} x_{2}+6 x_{1} x_{3}-6 x_{2} x_{3}
\]

其秩为 2 。

\begin{itemize}

  \item[1．] 求参数 \(c\) 的值．

  \item[2．] 求正交变换化 \(f\left(x_{1}, x_{2}, x_{3}\right)\) 为标准形．

  \item[3．] 指出方程 \(f\left(x_{1}, x_{2}, x_{3}\right)=1\) 表示何种二次曲面．
\end{itemize}
\end{document}
